\section*{From the genealogical closure to the physical realization}

The thesis of this work attains four coordinated and irreducible closures:
genealogical, operatorial, matricial, and physical. It does not consist in adjoining a
set theory, a theory of operator algebras, or a quantum formalism
to a prior numerical construction. The four developments arise from
the same causal precedence: APP constructs the dual positional and
arithmetic realization; oriented TRIT distinguishes return, the torsional present, and opening;
TPK preserves value, sheet, orientation, residue, quotient, carry, boundary,
route, register, memory, survival, phase, and gauge. Only thereafter are
the compatible section, its readers, and its representations formed. The governing diagram
is, therefore,
\[
\begin{gathered}
 \mathrm{APP}
 \longrightarrow \mathrm{TRIT}_{\rm orientado}
 \longrightarrow \mathrm{TPK}
 \longrightarrow \mathcal X_\infty^{\rm cal},\\[2mm]
 \mathcal X_\infty^{\rm cal}
 \xrightarrow{\;\mathfrak R\;}\mathbb R,
 \qquad
 \mathcal X_\infty^{\rm cal}\times\mathcal C_{\rm exc}
 \xrightarrow{\;\mathfrak E\;}\mathcal E_{\rm exc},\\[2mm]
 C(X_m)\xrightarrow{\;\operatorname{diag}\;}M_{9^m}(\mathbb C)
 \xrightarrow{\;a\mapsto a\otimes I_9\;}
 \mathrm{UHF}(9^\infty)
 \xrightarrow{\;\pi_\tau(\,\cdot\,)''\;}R_{\rm hyp},\\[2mm]
 \mathsf P_{\rm CT}^{+}
 \xrightarrow{\;\mathfrak R_{\rm MD}^{\rm disc,+}\;}
 \mathbf{SpecRoute}\times\mathsf{Hilb}_{\mathbb R,J}
 \times\mathbf B\mathfrak M_D^+\times\mathbf B G_C.
\end{gathered}
\]
The four lines do not form a single linear chain. Real evaluation and exceptional projection are sister readers of the enriched section; the tracial branch passes through the cylindrical algebras and their unital inclusions; the physical realization begins with prepared routes endowed with the structures required by each component. This separation of domains prevents the von Neumann factor from being deduced directly from a section, the exceptional projection from being identified with an action on digits, or a physical magnitude from being presented as though it were a TPK state without a realization morphism.

\subsection*{Genealogical Closure: Extensionality, Infinity, and Decision}

The first closure determines what it means to preserve an object across all depths. Given an inverse system of finite spaces \((X_m,p_m)\), a section \(x=(x_m)_m\) is not characterized solely by the equality of its published values. The fiber-descent theorem establishes the conditions under which a finite reader descends to a quotient and distinguishes surjectivity, constancy on fibers, and uniqueness of the descended map. König's lemma guarantees the existence of a branch when the levels are finite, nonempty, and finitely branching; uniqueness further requires the complete character that selects a single compatible prolongation.

On a fixed section, the decreasing cylinders provide an exact
realization of the passage \(0\!-1\!-\!\infty\): the support contracts, the
height increases, the mass remains unitary, and the cylindrical measures
converge weak-* to \(\delta_x\). This result coordinates the finite Kronecker delta,
the conditional expectations between levels, and the Dirac delta as a point
measure, without identifying the latter with a differential operator. The
prefix foundation precludes infinite descending chains in each finite
presentation; coinduction controls compatible prolongation. The von Neumann
ordinals are lifted with a route-memory fiber, and cardinality appears as a
decategorification that does not retain the entire genealogy.

The same distinction requires cylindrical power, closed power, and
full power to be separated, as well as set-theoretic choice, natural choice,
algebraic choice, and effective choice. Encoding in ZF/ZFC preserves the
constructions and permits comparison with Extensionality, Foundation,
Power Set, and Choice; it does not make the real line or the set-theoretic universe
the origin of APP--TRIT--TPK\@. In the effective domain, decidability of each
finite window, existence of the limit, uniqueness of an instance, and the
impossibility of a uniform total decider are four distinct assertions.
The Gödel--Turing reduction affects only the fourth: it does not invalidate particular
selectors already constructed. Analogously, the Gödel--Cohen comparison separates
the extension \(M\mapsto M[G]\) from the internal enrichment of a state and
avoids conflating the crossing of depth-dense sets with genericity over a
model. Borel cylinders, comparison hierarchies, and the projective
minimax passage are introduced with their hypotheses of compactness,
continuity, and consistency, not as automatic consequences of inverse
notation.

This closure also has an informational formulation. If an update
\(U\) of the complete state is bijective, then
\(H(X\mid U(X))=0\), and the ideal logical Landauer term vanishes. If a
reader \(q\) preserves only the visible shadow, the discarded information is
localized in \(H(X\mid q(X))\). In the operatorial formulation, a
trace-preserving automorphism preserves entropy, whereas a trace-preserving
conditional expectation \(E_N:M\to N\) satisfies
\[
 S_\tau(E_Nh)-S_\tau(h)=D_\tau(h\Vert E_Nh)\geq0.
\]
The cost corresponds to forgetting the fiber, not to the reversible transport of the
enriched state; this identity does not assert null total physical dissipation nor
reduced computational complexity.

\subsection*{Operatorial Closure: From the Nonadic Unit to Continuous Dimension}

The second closure operationally realizes the evolutionary conservation of unity. For the complete system, one considers
\[
 A_m=M_{9^m}(\mathbb C),
 \qquad \iota_m(a)=a\otimes I_9,
 \qquad
 \tau_m(a)=9^{-m}\operatorname{Tr}(a).
\]
Identity \(\tau_{m+1}\circ\iota_m=\tau_m\) simultaneously preserves
unity and trace. The inductive limit is the UHF algebra of supernatural type
\(9^\infty=3^\infty\), and the weak closure of its tracial GNS representation is
the hyperfinite factor \(R_{\rm hyp}\) of type \(II_1\). Values
\(\operatorname{rank}(P)/9^m\) become dense, and the trace of projectors in
\(R_{\rm hyp}\) realizes the entire interval \([0,1]\): a continuous dimension
emerges from a tower of discrete inclusions preserving unity.

This construction is not to be conflated with the marked branch
\(j_{m,r}(a)=a\otimes p_r\). The isometries of an individual fiber generate in
norm the compact operators on the Hilbertian limit and, under weak closure, a
factor of type \(I_\infty\). Selection of one sector preserves the mark, but not the
unit of the tower; the simultaneous lifting of the nine fibers is what
produces the UHF--\(II_1\) branch. Nor may the UHF realization be identified
without proof with the crossed-product realization of the profinite clock: a
canonical conjugacy requires the action, measure, amenability,
ergodicity, and isotropy to be declared. The orthomodular logic of projectors, Boolean
contexts, normal states, and entropic normalization
\[
 S_\tau(9^m\rho)=S_{\rm vN}(\rho)-m\log9
\]
are derived within this operator construction and keep separate the
matrix state, its tracial density, and its statistical publication.

Topology, dimension, and composition do not collapse in the factor. The closure
transversal is recorded by means of
\[
 \mathfrak C_{\rm HMT}
 =\bigl(C(\Omega_\infty),R_{\rm hyp},\mathcal G_{\rm TPK}\bigr).
\]
The commutative spectrum preserves points and cylinders; the tracial factor preserves dimension and statistics; the groupoid preserves concatenation, admissible inversion, and the provenance of routes. This threefold architecture is formulated after the double projection and the exceptional corridor. It neither founds nor replaces them: Archimedean evaluation and the Paley--Witt--Golay--Leech--\(V^\natural\)--Monster--Moonshine chain remain the two coordinated readings of the same enriched antecedent.

\subsection*{Matrix Closure: Depth, External Size and Dilation}

The third closure introduces simultaneous control of three indices. If \(m\) measures genealogical depth, \(n\) the exterior size of a grid, and \(s\) the matrix level of its entries, the family \(\mathfrak M_{m,s}^{(n)}\) is bigraded by \((m,s)\) when \(n\) remains fixed and trigraded by \((m,n,s)\) when the exterior size also varies. In particular,
\[
 \mathcal K_{m,s}
 :=\operatorname{UCP}\!\bigl(C(X_m),M_s(\mathbb C)\bigr)
\]
has two independent operations: the genealogical restriction
\(\rho_m^s(\Phi)=\Phi\circ\jmath_m\) and the matrix compression
\[
 \operatorname{mc}_{\boldsymbol V}(\Phi_1,\ldots,\Phi_k)(f)
 =\sum_rV_r^*\Phi_r(f)V_r,
 \qquad \sum_rV_r^*V_r=I_s.
\]
Both commute exactly:
\[
 \rho_m^s\!\left(\operatorname{mc}_{\boldsymbol V}(\Phi_1,\ldots,\Phi_k)\right)
 =\operatorname{mc}_{\boldsymbol V}\!\left(
   \rho_m^{s_1}(\Phi_1),\ldots,\rho_m^{s_k}(\Phi_k)\right).
\]
The identity is functorial and distinguishes refinement of history of reduction from the observer.

Every depth-compatible family determines a unique UCP map
over \(C(X_\infty)\) and admits a common Stinespring dilation. In the
diagonal corridor, a single projective spectral measure \(P\) satisfies
\[
 E_x^{(m)}=V^*P([x]_m)V
\]
at every depth and on every cylinder. The minimal dilation is unique up to
unitary equivalence and exactly preserves what is seen by the reader; a direct sum
provides a faithful ambient realization that preserves all distinctions of the
algebra. On the UHF limit, a compatible family of UCP maps likewise admits a common
prolongation and dilation. The conditional expectation
\[
 \mathbb E_m=\operatorname{id}\otimes\tau_9:
 M_{9^{m+1}}(\mathbb C)\longrightarrow M_{9^m}(\mathbb C)
\]
preserves the trace and discards the final factor in a controlled manner; its Stinespring environment records the omitted digit. In general, no uniform finite auxiliary memory exists for all depths. Moreover, faithfulness of a representation does not by itself imply APP--TPK dynamical faithfulness: the latter requires an intertwiner compatible with route, monodromy, and composition.

The matrix theory is instantiated in verifiable objects. The qutrit quantum
magic squares with parameters \((9,3)\) and \((12,3)\) have, respectively,
81 and 144 effects, exactly unital row and column sums, and
declared simultaneous dilations; the same POVM systems produce quantum magic
cubes of orders nine and twelve. The order-twelve case is organized through
the fibration \(12\to4\times3\), with its gauge and its explicit
equivariance condition. By another route, the sum and product conditional expectations of APP
generate the algebra
\[
 C^*(P_{\Sigma,2},P_{\Pi,2})
 \cong\mathbb C^4\oplus M_2(\mathbb C)\oplus M_2(\mathbb C),
\]
and at every level \(s\) the central block realizes the free matrix interval
\([5I_s,9I_s]\). These constructions distinguish quantum Latin squares,
magic squares, projective measures, positive measures, and semiclassicality;
no identification is obtained merely because the objects share order
nine.

The common spectral antecedent of two readers is formulated only after
both have been constructed. For finite maps
\(R_m:X_m\to\mathcal A_m\) and \(Q_m:X_m\to\mathcal E_m\), the incidences
\[
 P_{a,e}=1_{\{R_m=a\}}1_{\{Q_m=e\}}
\]
refine the two diagonal spectral measures and recover their marginals. The
result does not prove joint injectivity or require one reader to be a
compression of the other. In the noncommutative regime, two mutually
unbiased qutrit PVM systems do not possess a sharp joint PVM; simultaneous realization requires
compatible noise or a larger environment. This difference in type preserves the
HMT double projection: the matrix theorem classifies realizations of readers
and does not rewrite their foundational antecedent.

\subsection*{Physical Closure: Common Domain, Two Orientations, and Localization}

The fourth closure constructs the domain from which a TPK route can simultaneously receive spectral, Hilbert-space, dimensional, and Cartan realizations. Let \(\mathsf D_{\rm sp}\), \(\mathsf D_Q\), \(\mathsf D_D\), and \(\mathsf D_C\) be the reference domains of those four representations, with faithful forgetful functors into the category \(\mathsf P_{\rm TPK}^{\rm enr}\) of prepared states and registered routes. The fiber product
\[
 \mathsf P_{\rm CT}^{+}
 :=\mathsf D_{\rm sp}
 \times_{\mathsf P_{\rm TPK}^{\rm enr}}\mathsf D_Q
 \times_{\mathsf P_{\rm TPK}^{\rm enr}}\mathsf D_D
 \times_{\mathsf P_{\rm TPK}^{\rm enr}}\mathsf D_C
\]
imposes that the four components have exactly the same antecedent.
Their canonical projections are functorial adapters and produce
\[
 \mathfrak R_{\rm MD}^{\rm disc,+}
 =\bigl(\mathfrak S_+,\mathcal Q_+,\mathfrak D_+,
        \rho_{C,+}^{\rm disc}\bigr),
\]
with values, respectively, in spectral routes, real Hilbert spaces
with complex structure, the positive dimensional monoid, and the discrete
Poincare group. The construction asserts the functor on the common domain; it does not
assert that every TPK route automatically admits all four lifts.

On a path \(\gamma=e_1\cdots e_r\), the total variation
\[
 V_+(\gamma)=\sum_jw(e_j)
\]
and the net oriented transport
\[
 V_{\rm or}(\gamma)=\sum_j\varepsilon(e_j)w(\underline e_j)
\]
are sibling readers. The loop \(ee^\vee\) proves that the former does not,
in general, factor through the latter: it has zero oriented transport and
positive variation. To obtain an invertible realization, the domain is restricted
to those edges whose spectral and Hilbertian components are invertible, whose
dimensional datum is an additive cochain, and whose Cartan cocycles are normalized
and antisymmetric. Only then does the localization
\(\Pi_\vee(\mathsf P_{{\rm CT},{\rm rev}}^+)\) produce the oriented functor
\(\mathfrak R_{\rm MD}^{\rm disc,or}\) and transport inversion component by
component. The positive reading is not eliminated, and the oriented reading is not
obtained merely by changing a sign.

This physical construction precedes the Planck and electron constructions
because it supplies them with a domain, representations, and transport rules, but
does not absorb their reference proofs or the General Mass Law. The
number--particle--route--enriched chronological record--signature--
character--towers chain preserves its causal direction; the absolute action scale,
the central electron, the mass families, CKM, PMNS, and the physical sectors
retain their own objects, domains, and conclusions. A metrological coincidence is
compared only after the internal output has been frozen. Extensions that still
require an intertwiner, a smooth limit, a proof of dynamical stability,
or a physical morphism are stated with those data as exact hypotheses. The
absence of any of those maps prevents the theorem from being extended to the
corresponding codomain, without altering the genealogical, operatorial, and
matricial theorems already proved in their domains.

The reading of the work follows, consequently, the order of dependence and not that of external recognition:
\[
\begin{aligned}
 \mathrm{APP}&\longrightarrow\mathrm{TRIT}_{\rm orientado}
 \longrightarrow\mathrm{TPK}\\
 &\longrightarrow\text{enriched section}
 \longrightarrow
 \begin{cases}
  \text{Archimedean evaluation},\\
  \text{exceptional projection},\\
  \text{cylinder algebras and tracial closure},\\
  \text{prepared routes and physical realization}.
 \end{cases}
\end{aligned}
\]
Each output preserves its domain, hypotheses, and precise obstructions. What is common is not
a forced equality between codomains, but provenance from the same
enriched state. This is the precise meaning of holographic closure: no
terminal publication exhausts the genealogy that produces it, and no subsequent
closure can replace the double direction --Archimedean and exceptional--
that preserves the unit with memory.
